\section{Introduction}\label{sec:intro}
% Outline row: introduction.tex. ISS care is ground-directed; Mars light time; NASA's distance hazard and EIMO;
% alarm fatigue; the gap; four contributions in prose.

On the International Space Station, medical care is directed from the ground. Astronauts rely on real-time contact with medical providers in mission control, and most medical expertise and decision-making sit on Earth~\cite{russell2023cdss,sirek2021autonomy}. On board, care is given by the crew medical officer (CMO), a crew member for whom there is little time for medical training among everything else a mission demands~\cite{russell2023cdss}. The arrangement works because the delay is short.

At Mars it is not. The one-way light time alone is about 3 to 22 minutes, from the closest and farthest Earth--Mars distances~\cite{nssdcMars}, and across planned crewed trajectories it peaks at about 22 minutes~\cite{mcbrayer2022delay}. On 14 November 2026, the date of our demonstration mission, it is about 11 minutes (\numMarsOneWayS~s), and a question to Earth costs the round trip, twice that. NASA lists distance from Earth as one of the hazards of human spaceflight and expects crews bound for Mars to confront medical emergencies with minimal support from teams on Earth~\cite{hrpDistance}. Its plans for Earth-independent medical operations (EIMO) move onboard care progressively into the crew's hands as missions go further out~\cite{levin2023eimo,lemery2024eimo}.

Monitoring software will be part of that shift, but more alarms are not the answer. Bedside monitors in intensive care already raise so many alarms that staff suffer alarm fatigue, and a life-threatening event can be missed among them; in one large study most of the arrhythmia alarms that nurses annotated were false positives~\cite{drew2014alarm}. A crew of a few people with one CMO cannot absorb that. A monitor for deep space has to raise fewer alarms, and each one has to reach the person who can act on it in time.

What distance changes is who can act in time. NASA's NEEMO~16 analog mission ran a simulated medical emergency under a five-minute communication delay, and the crew member's safety was compromised by the delay between crew and flight surgeon~\cite{rader2013delay}. If a reading is falling and will reach a dangerous level within a couple of minutes, a flight surgeon whose answer is a round trip away, \numMarsRoundTripS~s at Mars on our date, cannot help, however good the data they receive. Frameworks for space medicine already describe autonomy that grows with distance and delay~\cite{levin2023eimo,cypko2025}. We did not find an open monitoring system that takes the light time as an input on every alert, compares it with how long the patient has, and says plainly who must act.

LIGHT-TIME, built for the NASA Space Apps Challenge 2026, is a small attempt at that, and this paper describes its decision layer and how we checked it. First, the ship computer sets who acts, Earth or the crew, from the one-way light time read from an ephemeris for the mission date; the switch point is a design choice, kept as one constant. Second, a time-left estimate built from a Theil--Sen slope and an agreement gate is compared with the round trip, so a warning says whether Earth can answer before the reading reaches its red line. Third, trust rules discard readings that cannot be believed before any band or trend is judged. Fourth, we replayed public physiological recordings, all from volunteers on Earth, through the same rules, report every event, and check the settings by leaving each recording out; the code, its tests and the scripts that produce every number in this paper are open. The evidence is small, and so is the claim: this is an engineering design with a first replay check, not a clinical validation.
